% !Mode:: "TeX:UTF-8"

\section{论文工作是否按开题报告预定的内容及进度安排进行}

论文工作总体按开题报告预定的内容与进度推进，无方向性调整。开题
设定的四项研究内容中：（1）目的地内转移行为统计建模已完成主体，
处于参数区间估计补充阶段；（2）融合方法研究超额完成，新增了开题
未明确列出的组件必要性消融与概率校准实验；（3）三层统计评测体系
已从设计文档落地为可复现实验代码并完成生成层验证；（4）系统接入
完成（规划器新增先验与融合两档选择器），端到端演示待论文收尾
阶段完善。进度上，开题安排的"2026.09--10 完成混合模型参数区间
估计与生成层扩展"中，生成层实验已提前完成，参数区间估计顺延至
中期检查后，总体进度正常略超前。

\section{目前已完成的研究工作及结果}

\subsection{行为统计建模（对应研究内容 1）}

以拟合份 84 条真实路线估计全部行为模型组件，测试份 42 条仅用于
最终评估。主要结果：

\begin{enumerate}
  \item \textbf{双峰性检验}：Hartigan dip test\cite{hartigan1985dip}
        在原始与对数尺度均拒绝单峰（$p<10^{-4}$）；路线级自助
        100/100 次拒绝\cite{efron1993bootstrap}。两分量对数正态
        混合未获似然比支持（$p=0.219$），提示跨区分量重于对数
        正态，最终采用 lognormal（中位 0.58km）与 log-gamma
        （中位约 25km）的异质混合，EM 估计\cite{dempster1977em}，
        就近分量权重 0.45。
  \item \textbf{转移阶数}：参数自助法似然比检验支持二阶区域转移
        （$p=0.005$）；样本外对数损失二阶 0.62 对一阶 0.75
        nat/转移。需要如实说明：在 top-1 命中率上二阶与一阶持平
        （37.1\% 对 36.8\%，配对 McNemar $p=1.0$），二阶的优势
        体现在似然质量而非命中率。
  \item \textbf{距离衰减选型}：候选池离散选择任务上的条件 logit
        拟合，混合形式优于幂律与指数（确认份对数损失 1.88 对
        2.08/2.09）；纯指数形式的极大似然估计退化为近平坦
        （尺度参数趋于上界），工程常用的固定指数衰减被数据否决。
  \item \textbf{区域粒度}：转移计数覆盖分析确定 8 区域（更粗粒度
        命中率下降约 4 个百分点）。
\end{enumerate}

\subsection{融合方法与架构必要性（对应研究内容 2）}

测试份 372 个预测点、8 选 1 任务，配对精确
McNemar\cite{mcnemar1947note}：

\begin{enumerate}
  \item \textbf{主结果}：统计先验单独命中率 37.1\%，显著高于最强
        规则基线（就近贪心）的 32.3\%（$p=10^{-4}$）；与 LLM
        候选概率对数池化后 43.3\%（对规则 $p\approx0$，对先验
        单独 $p=0.005$）。LLM 单独仅 26.6\%——其边际价值在融合
        框架下实现反转。
  \item \textbf{架构必要性消融}：距离衰减是命脉（去除后降至
        18.8\%）；区域项贡献 4.8 个百分点；区域数从 8 降到 4/6
        损失 4 个百分点；平滑常数在 0.05--0.5 内不敏感。衰减
        形式对先验单独影响小，但对融合影响大（混合形式 43.3\%
        对退化的指数形式 34.4\%）。
  \item \textbf{负结果（如实报告）}：状态依赖权重 $\lambda(s)$
        （逻辑回归堆叠 + 嵌套交叉验证）不优于固定权重（43.0\%
        对 43.3\%，$p=1.0$），学得权重塌缩为 $[0.52,0.57]$ 窄带；
        LLM 概率温度校准在测试集无实质增益（期望校准误差原始仅
        0.041，本任务非严重过度自信）。最终部署架构据此精简为
        四组件：二阶区域转移、混合距离衰减、LLM 首 token 候选
        概率、固定权重的对数池化。
\end{enumerate}

\subsection{三层统计评测体系（对应研究内容 3）}

按构念效度框架\cite{cronbach1955construct}将评测构念明确为
"行为符合度"，与"工程质量"分离；三层中两层不依赖任何模型
（裁决层），一层依赖行为模型设定（描述性层），三层互证控制循环
论证——该分层借鉴生成模型评测中似然与样本分布判据不可互替的
论证\cite{theis2016note}。生成层实验（60 条指令、5 种选择器）：

\begin{enumerate}
  \item \textbf{个体层（描述性）}：路线行为对数似然（nat/转移），
        真实测试路线参照 $-3.84$ 为最高（合理性自检通过）；
        先验选择器 $-4.02$、融合 $-4.14$、就近规则 $-4.59$、
        LLM $-4.79$。配对 Wilcoxon：先验与融合均显著高于就近
        规则（$p\approx0$）。
  \item \textbf{群体层（裁决）}：转移距离分布对真实参照的 KS
        距离：融合 0.219（最优，bootstrap 95\%CI $[0.210,
        0.248]$）$<$ 先验 0.248 $<$ 就近 0.253 $<$ LLM 0.270
        $<$ 检索序 0.300；真实对真实参照基线为 0.046。跨区占比
        （真实 0.49）：就近规则 0.50 形式上最接近，但其分布
        90 分位 4.4km 相对真实约 8km 存在系统性的重尾压缩——
        这是全部生成方法的共同局限（源于候选池以当前中心检索），
        论文将如实报告并列为改进方向。
  \item \textbf{旧工程指标（附录，描述性）}：在已知偏向就近基线
        的综合评分下，先验（0.78）与融合（0.75）仍高于就近规则
        （0.65）与 LLM（0.65），结论跨指标族稳健。
\end{enumerate}

\subsection{系统接入（对应研究内容 4）}

行程规划器新增"先验"与"融合"两档候选选择器，逐日生成时逐步
以对数转移概率、距离衰减与类型转移打分（融合档再与 LLM 首 token
概率池化，单次前向、确定性）。全部实验脚本与结果数据已版本化，
可复现。

\section{后期拟完成的研究工作及进度安排}

\begin{enumerate}
  \item 2027.01：混合模型参数区间估计与分量数选择的稳健性补充；
        生成层实验的多种子复验。
  \item 2027.02：任务适配微调敏感性实验（以融合增益为裁决指标，
        回应"LLM 基线可能被削弱"的质疑）；多重比较校正（Holm）
        固化到全部报告。
  \item 2027.03：Web 演示系统集成与案例分析；全部消融链复跑。
  \item 2027.03--2027.05：论文撰写、预答辩、修改与正式答辩。
\end{enumerate}

\section{存在的困难与问题}

\begin{enumerate}
  \item 样本量上限：真实路线 168 条，更细的分层检验功效不足，
        目前依靠功效前置设计与路线级自助缓解；
  \item 重尾压缩：生成路线缺失真实数据中的极端长转移，需改进
        候选检索的空间覆盖（拟在中期后实验全城候选池方案）；
  \item 时间风险：论文撰写与实验收尾并行，需严格控制微调敏感性
        实验的规模（单卡算力约束）。
\end{enumerate}

\section{如期完成全部论文工作的可能性}

核心建模、主实验与评测体系均已完成且结果显著，剩余工作为稳健性
补充、单一敏感性实验、系统集成与论文撰写，无方向性风险。评估
可以如期完成全部论文工作。
